\reportchapter{I. GIỚI THIỆU CHUNG}
\reportsection{a. Giới thiệu Công ty TNHH JITS Innovation Labs}
Công ty TNHH JITS Innovation Labs là doanh nghiệp hoạt động trong lĩnh vực phát triển phần mềm và cung cấp dịch vụ công nghệ thông tin tại Hà Nội. Công ty bắt đầu hoạt động từ tháng 6 năm 2016, có quy mô từ 11 đến 50 nhân sự, đặt trụ sở tại Hà Nội. Quy mô này tạo điều kiện để các nhóm kỹ thuật phối hợp trực tiếp, đồng thời yêu cầu mỗi thành viên chủ động trong quá trình nghiên cứu, phát triển và triển khai sản phẩm.

Lĩnh vực chuyên môn của JITS Innovation Labs là các giải pháp công nghệ cho ngân hàng và tổ chức tài chính, bao gồm nền tảng ngân hàng số, hệ thống điểm bán hàng, xử lý dữ liệu theo luồng và tích hợp hệ thống. Đội ngũ của công ty có kinh nghiệm triển khai, tùy biến hệ thống ngân hàng lõi cho khách hàng doanh nghiệp, hỗ trợ các ngân hàng và tổ chức tài chính vi mô cải tiến sản phẩm và thực hiện chuyển đổi số. Đặc thù của các hệ thống tài chính đòi hỏi phần mềm có tính ổn định, khả năng mở rộng, bảo mật và giám sát tốt; vì vậy, các phương pháp DevOps và DevSecOps có vai trò thiết thực trong hoạt động phát triển và vận hành sản phẩm.

Bên cạnh hoạt động phát triển sản phẩm, doanh nghiệp tổ chức chương trình thực tập theo hướng đào tạo gắn với thực hành. Tham gia vào kỳ thực tập năm nay, sinh viên chúng em được tiếp cận kiến thức theo lộ trình, hoàn thành các bài tập nhỏ cho từng chủ đề, lưu trữ sản phẩm trên hệ thống quản lý mã nguồn và trình bày kết quả định kỳ. Hình thức này giúp chúng em hình thành thói quen tự học, viết tài liệu, kiểm chứng kết quả và chịu trách nhiệm đối với sản phẩm đã thực hiện. Trong kỳ thực tập, em tham gia chương trình đào tạo DevOps và lựa chọn hướng chuyên sâu MLOps và Security để vận dụng kiến thức vào dự án tổng hợp.

Thông tin doanh nghiệp được đối chiếu từ trang chủ, trang LinkedIn và thông tin đăng ký doanh nghiệp công khai \cite{jits-home,jits-linkedin,business-registry}.

\reportsection{b. Giới thiệu công việc thực tập}
Trong thời gian thực tập tại JITS Innovation Labs, em đảm nhiệm vị trí DevOps Intern. Em vận dụng kiến thức về Linux, Git, Docker, CI/CD, Kubernetes, Terraform, MLOps và DevSecOps để hoàn thành các bài thực hành và dự án tổng hợp. Quá trình này giúp em hình thành cách tiếp cận công việc từ tổ chức mã nguồn, xây dựng môi trường chạy đến kiểm tra và vận hành sản phẩm.

Từ các kiến thức nền tảng, em xây dựng và đóng gói các ứng dụng web nhiều thành phần, triển khai trên Kubernetes, quản lý hạ tầng bằng Terraform theo mô hình Infrastructure as Code (IaC), kết hợp với kiểm thử tự động và cơ chế CI/CD. Em cũng ứng dụng thêm MLOps và bảo mật. vào việc pipeline xử lý dữ liệu, huấn luyện, đánh giá và triển khai mô hình.

Những kinh nghiệm trên được em ứng dụng để phát triển dự án cá nhân Philobiblus, một ứng dụng web hỗ trợ quản lý tiến độ đọc, nhận xét, chia sẻ và gợi ý sách. Sản phẩm kết hợp các nội dung đã thực hành để tổ chức mã nguồn, đóng gói ứng dụng, tự động kiểm thử và triển khai thử nghiệm. Qua đó, em có thể hoàn thiện một sản phẩm từ giai đoạn phát triển đến môi trường vận hành và đánh giá được các yêu cầu về tính ổn định, khả năng mở rộng và bảo mật.

\reportsection{c. Giới thiệu bài toán}

Trọng tâm của bài toán là kiểm chứng cách tổ chức và vận hành một ứng dụng web nhiều thành phần: hệ thống phải ổn định khi frontend, backend, cơ sở dữ liệu và model học máy hoạt động đồng thời; bảo vệ dữ liệu theo quyền sở hữu và phạm vi truy cập; cô lập và khôi phục thành phần gặp lỗi; có khả năng autoscale khi tải và số người dùng đồng thời biến động. 
Philobiblus là một hệ thống mẫu quản lý thư viện và hoạt động đọc, gồm cập nhật tiến độ, nhận xét, chia sẻ và gợi ý sách, được em phát triển như là một sản phẩm mẫu để kiểm chứng các yêu cầu trên. Hệ thống được triển khai trên Kubernetes, quản lý hạ tầng bằng Terraform, kết hợp với CI/CD, MLOps và DevSecOps.
